\documentclass[11pt]{article}
\usepackage{amsmath,amssymb,amsthm}
\newtheorem{theorem}{Theorem}
\title{Gluing reverse modules at different deletion depths}
\author{Joint Cell Head Compression}
\date{24 November 2026}
\begin{document}
\maketitle

\paragraph{Status.}
The transfer and singleton repair below are kernel checked in
\texttt{GradedShellTransfer.lean}. They are conditional compiler interfaces,
not a new asymptotic bound. The definition of the original model and of
\texttt{HeadComplexity.worstCaseHeadComplexity} is unchanged.

Identify an input with a subset of $[n]$. For a center $C$ and depth $d$ put
\[
 \mathcal S(C,d)=\{X\subseteq C:|C|=|X|+d\},
 \qquad \mathcal D(C)=\{X:X\subseteq C\}.
\]
A reverse module of depth $d$ is a sum of $H$ affine fractions whose
individual denominator slopes are strictly negative, whose denominators
are positive on the entire Boolean cube, and whose sum vanishes at every
input of weight less than $d$. Numerators may be arbitrary affine forms.
The structure \texttt{ReverseGradedData} records exactly these conditions.

\begin{theorem}
Let $(C_i,d_i,L_i)_{i<N}$ be a finite family of reverse modules, using
$H_i$ heads. Suppose $|C_i|-d_i$ is nonincreasing in $i$, the shells
$\mathcal S(C_i,d_i)$ cover the cube, and
\[
 \operatorname{sgn}(f(X))L_i(C_i\setminus X)>0
 \quad\text{for every }X\in\mathcal S(C_i,d_i).
\]
Then the original model computes $f$ with at most $\sum_{i<N}H_i$ heads.
If such modules exist for every independently specified Boolean truth
table $f$, then $W(n)\leq\sum_{i<N}H_i$.
\end{theorem}
\begin{proof}
Reflect the local coordinates in $C_i$. A reverse denominator becomes
an affine denominator with strictly positive slopes on those coordinates;
assign slope one on coordinates outside $C_i$. Positivity at zero follows
from the assumed full-cube positivity of the local denominator. Every
resulting denominator therefore remains positive on the full cube. Add a
sufficiently large multiple of the nonnegative affine penalty
$\sum_{j\notin C_i}x_j$ to suppress this module away from its downset.
Inside the downset, the reflected score is exactly $L_i(C_i\setminus X)$.

If an earlier target $X\in\mathcal S(C_j,d_j)$ lies in the later downset,
then the ordering gives $|X|\geq|C_i|-d_i$. Hence either
$|C_i\setminus X|<d_i$, where the new score is exactly zero, or equality
holds, in which case $X$ is also a target of the new module and its score
has the same prescribed strict sign. Successive positive scaling and
outside-downset suppression preserve every prior strict sign. The finite
strict-sign gluing theorem supplies finite parameters for all modules.
Thus no limiting score is substituted for an actual model score.
Overlaps need no disjointness assumption.
\end{proof}

The exact declarations are
\texttt{glue\_reverse\_graded\_shell\_modules},
\texttt{HStar\_le\_graded\_shell\_cover}, and
\texttt{worstCaseHeadComplexity\_le\_graded\_shell\_cover}, in namespace
\texttt{HeadComplexity.AxisZeroGluing}. The depth-zero compiler
\texttt{exists\_reverseGradedData\_depth\_zero} uses the single fraction
\[
 \frac{\operatorname{sgn}(\ell)}{n+1-\sum_jx_j}.
\]
Its denominator is at least one everywhere. It can repair a singleton
exception at cost one within the same rank ordering, without a separate
unproved claim that prior signs survive an exception repair.

\paragraph{What a sparse query family still needs.}
A constant-weight embedding by itself is insufficient. If a proposed
module decodes only $Q_i\subsetneq\mathcal S(C_i,d_i)$, an earlier query
may lie in $\mathcal S(C_i,d_i)\setminus Q_i$. Neither lower-weight
cancellation nor the local sign hypothesis controls that input. One must
supply signs on the remainder, prove its score zero there, or establish a
stronger geometric cover compatibility condition. Moreover, a denominator
gate that selects an upper coordinate face generally remains nonzero after
some inactive coordinates are deleted. It does not automatically provide
the required lower-weight cancellation. These are the two separate missing
conditions for the gated butterfly query construction.

\end{document}
